\section{Introduction}
\label{sec:intro}

An AI agent is not a single artifact. What it does is determined by a system prompt,
a policy bundle, a set of tool schemas, a model and its version, a retrieval corpus,
a memory state, the decisions it records, the sub-agents it delegates to, its supply
chain, and the human approvals that gate its actions. Today none of these is
attested. A signed JWT proves which credential called an API; it proves nothing about
the agent behind the call.

This is the \emph{agent identity gap}. When an agent acts in a regulated environment,
its operator can show policy documents, SOC reports, and mutable application logs,
none of which prove what actually ran. A reviewer cannot confirm which system prompt
defined the agent, which model version executed, which policy was in force, which
tools were authorized, whether the retrieval corpus or memory was poisoned, what the
agent decided, whether a human approved a high-risk action, or whether the deployed
binary matches the one reviewed.

Agent Manifest closes this gap. It binds the artifacts that define an agent into one
cryptographically signed, hardware-attestable, tamper-evident record. Its central
design choice is \emph{tamper-evidence over tamper-resistance}: rather than trying to
prevent change, the manifest makes change \emph{detectable} by any third party that
holds it and compares it to hardware attestation. Each artifact is bound by hash;
the manifest is canonicalized with RFC~8785~\citep{jcs8785}, signed with
Ed25519~\citep{eddsa8032} (or post-quantum ML-DSA-65~\citep{mldsa204}), and registered
on an append-only transparency log.

\subsection*{The ten unattested surfaces}

Agent Manifest binds ten artifacts: (1)~the system prompt, (2)~the policy bundle,
(3)~the tool manifest, (4)~the model identity, (5)~the retrieval (RAG) corpus,
(6)~the memory baseline, (7)~the decision trace, (8)~the A2A delegation chain,
(9)~the supply chain, and (10)~the human-in-the-loop approvals. Each is a surface an
adversary can alter without detection in today's deployments; each becomes
tamper-evident once bound (Section~\ref{sec:model}).

\subsection*{Contributions}

\begin{itemize}[leftmargin=1.5em]
  \item \textbf{A manifest schema} binding ten agent-defining artifacts, with a
    SPIFFE identity and a confirmation key (Section~\ref{sec:model}).
  \item \textbf{Cryptographic protocols}: RFC~8785 canonicalization, Ed25519 and a
    post-quantum ML-DSA-65 profile, and RFC~9162 Merkle commitments
    (Section~\ref{sec:crypto}).
  \item \textbf{A fail-closed verification protocol} with explicit result states and
    revocation (Section~\ref{sec:verification}).
  \item \textbf{Delegation and human-oversight} bindings that prevent scope laundering
    across agent hops (Section~\ref{sec:deleg}).
  \item \textbf{Integration and conformance}: composition with cMCP, TRACE, and
    AGT (Section~\ref{sec:integration}), and an open SDK with a conformance specification
    of 197 tests across four levels (Section~\ref{sec:conformance}).
\end{itemize}

\subsection*{Paper structure}

Section~\ref{sec:background} reviews the primitives Agent Manifest composes.
Section~\ref{sec:model} defines the manifest. Section~\ref{sec:crypto} gives the
cryptographic protocols and Section~\ref{sec:verification} the verification protocol.
Section~\ref{sec:deleg} covers delegation and human oversight,
Section~\ref{sec:integration} integration with the surrounding stack,
Section~\ref{sec:conformance} conformance, Section~\ref{sec:related} related work and
regulatory mapping, and Section~\ref{sec:conclusion} concludes.
